\section{更新过程与长期更新率}
\label{sec:06-Random-Processes/02-Counting-Processes/04}

机房里几千块硬盘，坏一块换一块。你想知道明年要备多少块。

上一节的危险率图已经排除了指数：硬盘的失效率随役龄上升，画出来是那条红线。间隔既然不是指数，无记忆性没了，Poisson 那一整套闭式解也就跟着没了。可备件预算这个问题并不复杂，凭什么它非要依赖那么强的假设？

答案是不依赖。只要每次更换之后系统在统计意义上从头开始，长期的更换频率就被平均寿命一个数字决定，间隔具体服从什么分布无关紧要。这是本节的主结果，也差不多是这个框架下唯一还能免费拿到的东西。

\subsection{只保留一条假设}

设相邻间隔 \(W_1,W_2,\ldots\) 独立同分布、取正值，分布任意。令

\[
T_n=W_1+\cdots+W_n,
\qquad
N(t)=\max\{n:T_n\le t\},
\]

\(\{N(t)\}\) 就叫更新过程。Poisson 过程是 \(W_i\sim\Exp(\lambda)\) 的特例。

先看丢了什么。若 \(W\) 带记忆，知道``当前这块盘已经转了多久''就会改变对剩余寿命的判断，于是 \(N(t)\) 单独不再构成一个足以预测未来的状态——它没记住当前周期的进度。要恢复 Markov 性，得把年龄

\[
A(t)=t-T_{N(t)}
\]

或剩余寿命 \(R(t)=T_{N(t)+1}-t\) 一并放进状态。这正是上一节说的``无记忆性让状态空间塌缩''反过来的样子：假设放松了，状态就得变胖。

独立增量也一起失守。相邻两个观察窗口会共享同一个跨越边界的周期，那个周期的长度同时影响两边的计数，两个窗口因此相关。平稳增量在有限时间内同样不成立——从一个更新点出发观察，起点附近的行为和很久以后的行为不一样。剩下的结构少得可怜：只有``每个更新点处重新开始''这一条。

\subsection{长期更新率只看均值}

设 \(\mu=\E[W_1]<\infty\)。由\hyperref[sec:05-Probability/08-Limit-Theorems/02]{``大数定律与长期平均''}，\(T_n/n\to\mu\)。要把它翻成关于 \(N(t)\) 的结论，用一次夹逼：按定义 \(T_{N(t)}\le t<T_{N(t)+1}\)，两边同除 \(N(t)\)，

\[
\frac{T_{N(t)}}{N(t)}\ \le\ \frac{t}{N(t)}\ <\ \frac{T_{N(t)+1}}{N(t)+1}\cdot\frac{N(t)+1}{N(t)}.
\]

\(t\to\infty\) 时 \(N(t)\to\infty\)，左右两端都收敛到 \(\mu\)，中间被夹住。取倒数得

\[
\frac{N(t)}{t}\ \longrightarrow\ \frac1\mu
\qquad\text{几乎必然}.
\]

论证只用到大数定律和那两个不等式，没有碰间隔的分布。平均每个周期耗时 \(\mu\)，长期每单位时间就完成 \(1/\mu\) 次更新——写出来近乎废话，但它对任何寿命分布都成立，这才是它有用的地方。硬盘的备件预算只需要平均寿命，不需要知道失效曲线的形状。

另有一条形式相近的结论，\(\E[N(t)]/t\to1/\mu\)，通常称作基本更新定理。两者极限相同，含义不同：前者说的是几乎每一条样本路径的长期频率，后者说的是期望值的增长速度。而且不能由前者直接取期望推出后者——几乎必然收敛不自动给出期望的收敛，中间需要一致可积之类的条件。这类地方是应用中最容易含混过去的：你在一台机器上观察到的是路径频率，你在预算表上写的是期望，两者相等需要额外的理由。

把报酬挂上去，结论会变得更好用。若第 \(i\) 个周期附带一个报酬 \(R_i\)，且 \((W_i,R_i)\) 作为数对独立同分布（同一周期内两者可以相关），则长期单位时间平均报酬为

\[
\frac{\E[R]}{\E[W]}.
\]

换盘成本 \(C\)、平均寿命 \(\mu\)，长期单位时间的更换开销就是 \(\E[C]/\mu\)；一个批处理作业平均跑 \(\mu\) 分钟、平均消耗 \(\E[R]\) 核时，长期的核时占用率同理。排队论里的 Little 定律也可以从这个角度理解——把``顾客在系统中停留的时间''当成周期报酬累积起来。一大类时间平均问题都能压缩成一句话：一个周期的平均报酬除以一个周期的平均长度。

\subsection{随机时刻会撞进长周期}

现在换一个问题，它的答案第一次看会觉得算错了。等到过程跑了很久，随手指一个时刻，你落在的那个周期有多长？

不是 \(\E[W]\)。长周期占据的时间更多，因此更容易被撞见。

\begin{center}
\begin{tikzpicture}[x=1cm,y=1cm,font=\small,>=stealth]
  \fill[blue!12] (1.0,-0.18) rectangle (3.4,0.18);
  \fill[blue!12] (4.3,-0.18) rectangle (7.6,0.18);
  \draw[->,black!70] (-0.3,0) -- (10.7,0) node[below,font=\footnotesize] {时间};
  \foreach \t in {0,0.6,1.0,3.4,3.9,4.3,7.6,8.2,10}{
    \draw[black!70] (\t,-0.22) -- (\t,0.22);
    \fill (\t,0) circle (1.8pt);}
  \draw[->,red!75,thick] (5.9,1.15) -- (5.9,0.30);
  \node[above,font=\footnotesize,red!75] at (5.9,1.16) {随手指一个时刻};
  \draw[black!60] (4.3,-0.58) -- (4.3,-0.75) -- (5.9,-0.75) -- (5.9,-0.58);
  \node[below,font=\footnotesize] at (5.1,-0.73) {年龄 \(A(t)\)};
  \draw[black!60] (5.9,-0.58) -- (5.9,-0.75) -- (7.6,-0.75) -- (7.6,-0.58);
  \node[below,font=\footnotesize] at (6.75,-0.73) {剩余 \(R(t)\)};
\end{tikzpicture}

\emph{八个周期，只有两个是长的（图中着色），按个数算占四分之一；但它们盖住了时间轴的 \(57\%\)。闭着眼睛在轴上指一个点，落进长周期的机会远超四分之一。你所在周期的年龄与剩余寿命把这个周期分成两段，两段都比``平均周期的一半''要长。}
\end{center}

写成分布就是长度偏置：被撞见的周期长度 \(L\) 的密度正比于 \(\ell f_W(\ell)\)，权重就是这个周期占了多长时间。归一化常数恰好是 \(\E[W]\)，于是

\[
\E[L]=\frac{\E[W^2]}{\E[W]}=\E[W]+\frac{\Var(W)}{\E[W]}\ \ge\ \E[W],
\]

等号当且仅当 \(W\) 是常数。多出来的那一项直接由方差控制：间隔越不齐，随机观察者所在的周期就越比典型周期长。在平衡状态下（需要 \(\E[W^2]<\infty\)、分布非格点），年龄与剩余寿命的平均都是 \(\E[W^2]/(2\E[W])\)，两者相加正好还原 \(\E[L]\)，不矛盾。

公交车平均十分钟一班，不代表你随便走到站台平均等五分钟。如果发车间隔时紧时松，你更可能撞进那个二十分钟的空档里，实际平均等待可以显著超过五分钟。同一件事在别处叫检查悖论：抽查正在运行的作业，抽到的多半是长作业；按时间点采样活跃会话，采到的多半是长会话；从当前在线用户里估计会话时长，估出来的数会系统性偏大。修正办法是按周期采样而不是按时刻采样，或者在估计时显式除掉 \(\ell\) 这个权重。

指数分布在这里是个有意思的特例。无记忆性让剩余寿命仍然服从原来的分布，\(\E[R]=\E[W]\)，看上去``没有偏置''；但年龄也服从同一分布，两段加起来是 \(2\E[W]\)——观察到的周期照样被拉长了一倍。所以偏置并没有消失，只是被无记忆性藏进了对称性里。

\subsection{连``更新''都不成立的时候}

整节只用了一条假设：周期独立同分布。它比无记忆性弱得多，但仍然可以被现实推翻。维修不彻底，修完之后的设备比新的更容易再坏，周期就不同分布了；环境随季节变化，寒冬里的失效间隔整体偏短，周期之间也不再独立；机器学习系统里的模型退化更麻烦，重训之后的表现与上一轮的数据分布有关，重训周期彼此纠缠。

对应的扩展各有名字：起点不在更新点的用延迟更新过程，运行与维修交替的用交替更新过程，周期类型随状态转移的用 Markov 更新过程。选哪个之前先问一句实质的：这次事件之后，系统是不是以与上次完全相同的统计状态重新开始？答不上来，就说明``更新''这个词还没有落实，套公式只会得到一个看起来精确的错数。

诊断上有个粗糙但管用的办法：把间隔序列按发生顺序排开，做一次自相关，再把它分成前后两半比较分布。有序列相关说明不独立，前后半分布明显不同说明不同分布。两项都过，独立同分布这条假设才算勉强站住。

本章到此为止：从最强的假设（Bernoulli 的固定 \(p\)）一路松到最弱的（只要求周期独立同分布），每松一步都换来更大的适用范围和更少的可算结论。松到底之后剩下的路只有一条——不再指望用一个分布描述整段历史，而是承认系统会在若干个状态之间迁移，把``现在处于哪个状态''显式记下来。这正是\hyperref[ch:048]{``离散时间 Markov 链：状态足够好，历史就可以被压缩''}的出发点。
